% ARCHIVED DRAFT: independently audited by root and incorporated 2026-11-19 18:01 UTC.
% Main part_B.tex is authoritative.
\begin{saresult}{B.080}{Proved under stated hypotheses}{Canonical arithmetic Tor obstructions for a flat effective coefficient ring}
Let $R\subset S$ be any unital commutative inclusion. Form the explicit
coefficient ring $C=\bigcup_d(R:_S D_d(R))$ of B.064, and assume that
$R\to C$ is flat and epic. Put $Q=C/R$ and define the $R$-module
\[
 E=\left\{\left(r\longmapsto\sum_{i=0}^d c_i r^i+R\right):
                       d\ge0,\ c_i\in C\right\}
       \subset\operatorname{App}(R,Q).                   \tag{B80.1}
\]
Thus $E$ is specified by coefficient and value relations in $R\subset C$,
without choosing tensor presentations of integer-valued polynomials.
For $A=\Omega(R,S)$ there is a canonical exact sequence
\[
\begin{gathered}
 0\longrightarrow\operatorname{Tor}_2^R(E,E)
 \longrightarrow A\otimes_R A\\
 \xrightarrow{\theta_2} C[X,Y]
 \xrightarrow{\partial}\operatorname{Tor}_1^R(E,E)
 \longrightarrow0.
\end{gathered}                                          \tag{B80.2}
\]
The map $\partial$ is given explicitly below, with all connecting maps
fixed by the evaluation sequence. In particular the two tensors on
153/004 exist uniquely exactly when
\[
\begin{gathered}
 \operatorname{Tor}_2^R(E,E)=0,\\
 \partial\!\left(\sum_{i,j\ge0}\binom{i+j}{i}
                         a_{i+j}X^iY^j\right)=0,\\
 \partial\!\left(\sum_{i\ge0}a_iX^iY^i\right)=0
                                                        \tag{B80.3}
\end{gathered}
\]
for every finite vector $(a_i)$ in $C$ satisfying
$\sum_i a_i r^i\in R$ for all $r\in R$; missing coefficients are zero.
If only flatness of $C$ is initially assumed, epicness is an additional
necessary condition, and epicness together with (B80.3) is sufficient.
This is an exact obstruction calculation for the flat-coefficient
scope, not a classification of arbitrary nonflat coefficient extensions.
\end{saresult}
\begin{proof}
Let $\mathcal F=C[X]$, and let
$\varepsilon:\mathcal F\to E$ send a coefficient vector to its
function modulo $R$. By the definition (B80.1),
\[
                0\longrightarrow A\longrightarrow\mathcal F
                    \xrightarrow{\varepsilon}E\longrightarrow0
                                                        \tag{B80.4}
\]
is exact. Also $E$ is a quotient of $Q[X]$, since every ordinary
$R$-coefficient polynomial maps to zero. Epicness gives
$C\otimes_R C\simeq C$, and the sequence $R\to C\to Q\to0$ then
gives $C\otimes_R Q=0$. Therefore
\[
 C\otimes_R E=0,\qquad \mathcal F\otimes_R E=0.           \tag{B80.5}
\]
Since $C$ is flat, so is $\mathcal F$. Tensoring (B80.4) with
$\mathcal F$ gives an isomorphism
\[
 \mathcal F\otimes_R A\ \xrightarrow{\sim}
 \mathcal F\otimes_R\mathcal F\simeq C[X,Y].             \tag{B80.6}
\]
Here the last map multiplies coefficients, using the ring epimorphism;
the variables in the two factors remain distinct. Explicitly the inverse
sends $H=\sum_j G_j(X)Y^j$ to $\sum_jG_j(X)\otimes Y^j$, since
$R[Y]\subset A$.

We use the long exact Tor sequence with its usual connecting maps.
For clarity, it is obtained by tensoring a free resolution with the
short exact sequence and applying the connecting homomorphisms of the
resulting degreewise exact sequence of complexes; this fixes all signs.
Flatness makes every positive Tor group with $\mathcal F$ vanish.
Applying this construction to (B80.4), first with the module $A$,
gives
\[
 0\longrightarrow\operatorname{Tor}_1^R(E,A)
 \longrightarrow A\otimes_R A
 \longrightarrow\mathcal F\otimes_R A
 \longrightarrow E\otimes_R A\longrightarrow0.          \tag{B80.7}
\]
Applying it with $E$ and using (B80.5) gives the canonical isomorphisms
\[
\begin{aligned}
 \operatorname{Tor}_2^R(E,E)&\xrightarrow{\sim}
                         \operatorname{Tor}_1^R(E,A),\\
 \tau:\operatorname{Tor}_1^R(E,E)&\xrightarrow{\sim}E\otimes_R A.
\end{aligned}                                           \tag{B80.8}
\]
Combining (B80.6)--(B80.8) gives (B80.2). These uses of Tor agree
with the opened Stacks Section 10.75, Lemmas 10.75.2, 10.75.5 and
10.75.8 \cite{StacksTor}; the construction above specifies the maps actually used.

The promised formula, containing only the displayed coefficient array
and the fixed connecting isomorphism, is
\[
 \partial(H)=\tau^{-1}
       \left(\sum_j\varepsilon(G_j(X))\otimes Y^j\right),
             \qquad H=\sum_jG_j(X)Y^j.                  \tag{B80.9}
\]
Every term $Y^j$ here is an ordinary monomial in the specified copy
$R[Y]\subset A$. Thus no unknown co-operation or decomposition of
$H$ into integer-valued factors enters the definition of $\partial$.
One can compute its class with any free resolution of the explicitly
defined function module $E$; the comparison maps between resolutions
identify the resulting classes. Every individual relation or boundary
in such a computation uses finitely many coefficients. This statement
does not claim finite generation of $E$, a uniform degree bound, or a
terminating decision algorithm for arbitrary presented rings.

The kernel in (B80.2) proves that uniqueness is equivalent to the
first condition in (B80.3). The other two conditions say, by the exact
sequence, exactly that the binomial coefficient array of $F(X+Y)$
and the diagonal coefficient array of $F(XY)$ lift. When the kernel
vanishes their lifts are unique, so (B80.9) recovers precisely the two
specified polynomial tensors. If epicness was not initially assumed,
B.064 proves it from uniqueness, giving the variant in the statement.
No higher tensor injectivity or plethory identities are inferred from
this binary calculation alone.

There is also an intrinsic distinction between the two specified
existence tests and full bivariate polynomial compositeness.
Let $\varepsilon_2:C[X,Y]\to\operatorname{App}(R^2,Q)$ be evaluation
modulo $R$. It vanishes on the image of $\theta_2$, so (B80.2)
defines a unique map
\[
 \lambda:\operatorname{Tor}_1^R(E,E)\longrightarrow
                  \operatorname{App}(R^2,Q),\qquad
                         \lambda\partial=\varepsilon_2. \tag{B80.10}
\]
Writing $\Omega_2=\ker\varepsilon_2$, exactness gives
\[
\begin{gathered}
 0\longrightarrow\operatorname{Tor}_2^R(E,E)
 \longrightarrow A\otimes_R A\\
 \longrightarrow\Omega_2
 \longrightarrow\ker\lambda\longrightarrow0.
\end{gathered}                                          \tag{B80.11}
\]
Indeed every class in $\ker\lambda$ has a preimage under the
surjective map $\partial$, and that preimage lies in $\Omega_2$.
Thus the full bivariate tensor map is an isomorphism exactly when
$\operatorname{Tor}_2^R(E,E)=0$ and $\lambda$ is injective. This
identification makes the remaining existence obstruction explicit
without assuming that the stronger condition follows from (B80.3).
The nonflat example B.067 is outside the flat-coefficient hypothesis,
so the present calculation does not close the original arbitrary-ring
structural task.
\end{proof}
